\scanpage{185}
\mspage{130}
% archive sheet number: 184

ment l'hypothèse $X\in\mathrm{Ob}\,\underline{\mathrm{Fib}}\,S$. Donc on ne peut espérer \unsure{donner|trouver} une description de $u^*$ \inserted{\deleted{\unsure{sur}} $\mathrm{HOT}_W(S)$ et $\mathrm{HOT}_W(S')$,} en termes d'opérations \deleted{\unsure{entre}} \inserted{sur les} catégories sur $S,S'$, \deleted{\unsure{quel}} comme suggéré à l'instant, qu'en se bornant à des ``modèles'' $X$ qui sont dans $\underline{\mathrm{Fib}}\,S$, \deleted{\unsure{tout}} au moins. Et je ne doute pas en effet que ça marche sur $\underline{\mathrm{Fib}\,S}$ pas moins bien que directement sur $\underline{\mathrm{Fibsc}}\,S=$ en termes de $\underline{\mathrm{Hom}}(S^{\unsure{\mathrm o|\circ}},\mathrm{Cat})$.

\centerline{\(\backslash\)}

Pour les axiomes essentiels sur les dérivateurs, je \unsure{renvois|renvoie} à I p. 67-68 et p. 111-112. On a vu que Der 1, \textcircled{\scriptsize Der 2}\textsuperscript{$*$} sont satisfaits par un prédérivateur défini par un couple $(\mathcal M,W\subset\mathrm{Fl}(\mathcal M))$, donc c'est OK ici, \inserted{pour $\mathrm{HOT}_W$}. Il s'ensuit que c'est OK aussi pour $\mathrm{HOT}^{\unsure{\mathrm{lc}|\mathrm{le}}}_W$, \deleted{qui} vu sa caractérisation \inserted{des $\mathrm{HOT}^{\mathrm{lc}}_W(S)$} en \deleted{\unsure{tous}} par sous-catégories des $\mathrm{HOT}_W(S)$. Pour Dér 1, ça marche \inserted{-il} même pour des sommes quelconques de catégories ?, i.e. a-t-on

% slanted
\margin{$*$ \unsure{Cavalier ça} pour \uline{Der} 2 \unsure{car}, je ne sais si c'est \unsure{tj.} vrai même en supposant $W$ \unsure{très} \uline{fortement} saturé \unsure{cf.}, i.e. tel que $\mathrm{hot}_W(f)$ iso $\Rightarrow f\in W$.}

\[
\mathrm{HOT}_W\bigl(\coprod_i S_i\bigr)\xrightarrow{\approx}\prod_i\mathrm{HOT}_W(S_i).
\]

On a en tous cas

\[
\underline{\mathrm{Fib}}\bigl(\overbrace{\coprod_i S_i}^{S}\bigr)
 \xrightarrow{\approx}\prod_i\mathrm{Fib}(S_i)
\]

et le localiseur $W^f_S$ \deleted{\unsure{en}} dans $\underline{\mathrm{Fib}}(S)$ s'identifie
